\documentclass[a4paper,12pt]{article}

% -------------------------------------------------
% БАЗОВАЯ КОДИРОВКА И МАТЕМАТИКА
% -------------------------------------------------
\usepackage[T2A]{fontenc}
\usepackage[utf8]{inputenc}
\usepackage[russian]{babel}

\usepackage{amsmath,amssymb,amsthm,mathtools}
\usepackage{icomma}

% -------------------------------------------------
% ТИПОГРАФИКА
% -------------------------------------------------
\usepackage{helvet}
\renewcommand{\familydefault}{\sfdefault}

\usepackage[
    a4paper,
    left=20mm,
    right=20mm,
    top=20mm,
    bottom=22mm
]{geometry}

\usepackage{microtype}
\usepackage{parskip}
\usepackage{setspace}
\onehalfspacing

\usepackage{enumitem}
\setlist[itemize]{
    leftmargin=6mm,
    itemsep=2pt,
    topsep=4pt
}
\setlist[enumerate]{
    leftmargin=8mm,
    itemsep=3pt,
    topsep=4pt
}

% -------------------------------------------------
% ГРАФИКА
% -------------------------------------------------
\usepackage{tikz}
\usetikzlibrary{
    calc,
    arrows.meta,
    positioning,
    shapes.geometric
}

\usepackage{tkz-euclide}

\usepackage{pgfplots}
\pgfplotsset{compat=1.18}

% -------------------------------------------------
% ТАБЛИЦЫ
% -------------------------------------------------
\usepackage{tabularx,booktabs,longtable}
\renewcommand{\arraystretch}{1.30}

% -------------------------------------------------
% СЛУЖЕБНОЕ ОФОРМЛЕНИЕ
% -------------------------------------------------
\usepackage{xcolor}
\usepackage{fancyhdr}
\usepackage{hyperref}

\definecolor{accent}{HTML}{008080}
\definecolor{darkaccent}{HTML}{004D4D}
\definecolor{textgray}{HTML}{333333}
\definecolor{softgray}{HTML}{6F7777}
\definecolor{lightaccent}{HTML}{DDEEEE}

\color{textgray}

\hypersetup{
    colorlinks=true,
    linkcolor=darkaccent,
    urlcolor=darkaccent
}

\pagestyle{fancy}
\fancyhf{}
\fancyhead[L]{\small\color{softgray}Графики функции и производной}
\fancyhead[R]{\small\color{softgray}Николь Саркисьянц}
\fancyfoot[C]{\small\color{softgray}\thepage}
\renewcommand{\headrulewidth}{0.3pt}
\renewcommand{\footrulewidth}{0pt}

\setlength{\headheight}{15pt}

% -------------------------------------------------
% ЗАГОЛОВКИ
% -------------------------------------------------
\newcommand{\maintitle}[1]{%
    {\Large\bfseries\color{darkaccent}#1\par}
    \vspace{2mm}
    {\color{accent}\rule{\linewidth}{0.6pt}}
    \vspace{4mm}
}

\newcommand{\subhead}[1]{%
    \vspace{3mm}
    {\large\bfseries\color{darkaccent}#1\par}
    \vspace{1mm}
}

\newcommand{\keypoint}[1]{%
    \par
    \vspace{1mm}
    {\color{darkaccent}\bfseries #1}
}

% -------------------------------------------------
% СТИЛИ БЛОК-СХЕМ
% -------------------------------------------------
\tikzset{
    flowstart/.style={
        ellipse,
        draw=darkaccent,
        line width=0.9pt,
        minimum width=39mm,
        minimum height=11mm,
        align=center,
        font=\small
    },
    flowaction/.style={
        rounded corners=2mm,
        rectangle,
        draw=accent,
        line width=0.9pt,
        minimum height=13mm,
        text width=56mm,
        inner sep=3mm,
        align=center,
        font=\small
    },
    flowdecision/.style={
        diamond,
        aspect=2.25,
        draw=darkaccent,
        line width=0.9pt,
        text width=35mm,
        inner sep=1.5mm,
        align=center,
        font=\small
    },
    flowfinal/.style={
        rounded corners=2mm,
        rectangle,
        draw=darkaccent,
        line width=1pt,
        minimum height=13mm,
        text width=68mm,
        inner sep=3mm,
        align=center,
        font=\small
    },
    flowarrow/.style={
        -{Stealth[length=3mm,width=2mm]},
        line width=0.85pt,
        draw=textgray
    }
}

% -------------------------------------------------
% ДОКУМЕНТ
% -------------------------------------------------
\begin{document}

% =================================================
% ТИТУЛЬНЫЙ ЛИСТ
% =================================================
\begin{titlepage}
\thispagestyle{empty}

\begin{center}

\vspace*{24mm}

{\color{accent}\rule{38mm}{1pt}}

\vspace{11mm}

{\Huge\bfseries\color{darkaccent}
Чек-лист
\par}

\vspace{6mm}

{\LARGE\bfseries
Графики функции и производной
\par}

\vspace{4mm}

{\Large
Знак производной, монотонность\\
и точки экстремума
\par}

\vspace{18mm}

{\large
Подготовка к ЕГЭ
\par}

\vfill

{\large
\today
\par}

\vspace{12mm}

{\large\bfseries
Лёвин Артём Александрович
\par}

\vspace{2mm}

{\large
эксклюзивно для Николь Саркисьянц
\par}

\vspace{25mm}

\end{center}

\begin{tikzpicture}[remember picture,overlay]
    \draw[
        accent!55,
        line width=1.2pt
    ]
    ($(current page.south west)+(18mm,18mm)$)
    .. controls
    ($(current page.south west)+(65mm,29mm)$)
    and
    ($(current page.south east)+(-72mm,7mm)$)
    ..
    ($(current page.south east)+(-18mm,18mm)$);
\end{tikzpicture}

\end{titlepage}

% =================================================
% ОСНОВНОЙ ЧЕК-ЛИСТ
% =================================================
\maintitle{1. Главная идея}

В задачах этого типа первое действие всегда одно:

\[
\boxed{\text{Определить, чей график дан: }f(x)\text{ или }f'(x).}
\]

От этого зависит вся последующая логика.

Связь функции и её производной:

\[
f'(x)>0
\quad\Longrightarrow\quad
f(x)\text{ возрастает},
\]

\[
f'(x)<0
\quad\Longrightarrow\quad
f(x)\text{ убывает}.
\]

Для дифференцируемой функции в рассматриваемой внутренней точке локального
экстремума выполняется необходимое условие

\[
f'(x_0)=0.
\]

При этом равенство

\[
f'(x_0)=0
\]

само по себе ещё не гарантирует наличие экстремума.

\subhead{Что означает знак производной}

\begin{center}
\begin{tabularx}{0.92\linewidth}{@{}>{\bfseries}l X@{}}
\toprule
Условие & Что происходит с функцией \\
\midrule
$f'(x)>0$ &
график функции идёт слева направо вверх; функция возрастает \\

$f'(x)<0$ &
график функции идёт слева направо вниз; функция убывает \\

$f'(x)=0$ &
касательная горизонтальна; точка может оказаться экстремумом или стационарной точкой другого типа \\
\bottomrule
\end{tabularx}
\end{center}

\keypoint{Важно.}
Знак производной сообщает направление изменения функции.
Точные значения производной по одному лишь факту возрастания или убывания
восстановить нельзя.

% -------------------------------------------------
\maintitle{2. Геометрический смысл производной}

Если функция дифференцируема в точке \(x_0\), то

\[
f'(x_0)
\]

равна угловому коэффициенту касательной к графику функции в точке

\[
\bigl(x_0,f(x_0)\bigr).
\]

Поэтому:

\[
f'(x_0)>0
\]

соответствует касательной, поднимающейся слева направо;

\[
f'(x_0)<0
\]

соответствует касательной, опускающейся слева направо;

\[
f'(x_0)=0
\]

соответствует горизонтальной касательной.

Чем больше модуль производной,

\[
|f'(x)|,
\]

тем больше по модулю наклон касательной.

\keypoint{Практический вывод.}
При построении качественного эскиза \(f'(x)\) по графику \(f(x)\)
сначала надёжно устанавливают знак производной.
Дополнительная информация о крутизне графика помогает уточнить форму
эскиза, если она действительно считывается с рисунка.

% =================================================
% ГРАФИКИ
% =================================================
\maintitle{3. Один пример: функция и её производная}

Рассмотрим функцию

\[
f(x)=x^3-3x.
\]

Её производная:

\[
f'(x)=3x^2-3.
\]

Корни производной:

\[
3x^2-3=0
\quad\Longleftrightarrow\quad
x=\pm1.
\]

\begin{center}
\begin{tikzpicture}
\begin{axis}[
    width=0.93\linewidth,
    height=7cm,
    axis lines=middle,
    xmin=-2.2,
    xmax=2.2,
    ymin=-3.2,
    ymax=3.2,
    xlabel={$x$},
    ylabel={$y$},
    samples=250,
    unit vector ratio=1 1,
    clip=false,
    xtick={-2,-1,0,1,2},
    ytick={-3,-2,-1,0,1,2,3},
    grid=major,
    grid style={line width=.2pt,draw=gray!20}
]
\addplot[
    domain=-2.05:2.05,
    line width=1.2pt,
    color=accent
]
{x^3-3*x};

\addplot[
    dashed,
    color=softgray,
    domain=-3.2:3.2
]
coordinates {(-1,-3.2) (-1,3.2)};

\addplot[
    dashed,
    color=softgray,
    domain=-3.2:3.2
]
coordinates {(1,-3.2) (1,3.2)};

\node[
    anchor=south west,
    text=darkaccent
] at (axis cs:1.25,1.7) {$y=f(x)$};
\end{axis}
\end{tikzpicture}
\end{center}

На графике видно:

\[
f(x)\text{ возрастает на }(-\infty,-1),
\]

\[
f(x)\text{ убывает на }(-1,1),
\]

\[
f(x)\text{ возрастает на }(1,+\infty).
\]

Следовательно,

\[
f'(x)>0
\quad\text{при}\quad
x<-1\ \text{или}\ x>1,
\]

\[
f'(x)<0
\quad\text{при}\quad
-1<x<1.
\]

\begin{center}
\begin{tikzpicture}
\begin{axis}[
    width=0.93\linewidth,
    height=6.7cm,
    axis lines=middle,
    xmin=-2.2,
    xmax=2.2,
    ymin=-3.6,
    ymax=9.5,
    xlabel={$x$},
    ylabel={$y$},
    samples=250,
    unit vector ratio=1 1,
    clip=false,
    xtick={-2,-1,0,1,2},
    ytick={-3,0,3,6,9},
    grid=major,
    grid style={line width=.2pt,draw=gray!20}
]
\addplot[
    domain=-2.05:2.05,
    line width=1.2pt,
    color=accent
]
{3*x^2-3};

\node[
    anchor=south west,
    text=darkaccent
] at (axis cs:1.15,4.8) {$y=f'(x)$};
\end{axis}
\end{tikzpicture}
\end{center}

% -------------------------------------------------
\maintitle{4. Если дан график функции \(f(x)\)}

Здесь мы смотрим прежде всего на \textbf{монотонность функции}.

\subhead{Основная схема рассуждения}

\[
\boxed{
\begin{array}{c}
f(x)\text{ возрастает}\\[1mm]
\Downarrow\\[1mm]
f'(x)>0
\end{array}}
\qquad
\boxed{
\begin{array}{c}
f(x)\text{ убывает}\\[1mm]
\Downarrow\\[1mm]
f'(x)<0
\end{array}}
\]

Если гладкий график переходит:

\[
\text{убывание}\longrightarrow\text{возрастание},
\]

то возникает локальный минимум, причём в стандартной дифференцируемой
ситуации

\[
f'(x_0)=0.
\]

Если график переходит:

\[
\text{возрастание}\longrightarrow\text{убывание},
\]

то возникает локальный максимум.

\subhead{При построении эскиза \(f'(x)\)}

Работайте в следующем порядке:

\begin{enumerate}
    \item Разбейте график \(f(x)\) на промежутки монотонности.
    \item На каждом промежутке определите:
    \[
    \text{возрастание или убывание}.
    \]
    \item Переведите это в знак производной.
    \item Отметьте точки, где достоверно известно, что
    \[
    f'(x)=0.
    \]
    \item Постройте качественный эскиз производной.
\end{enumerate}

\keypoint{Ограничение метода.}
По грубому эскизу функции обычно можно уверенно восстановить знак
производной. Форма графика \(f'(x)\) между характерными точками
может быть определена лишь качественно.

% -------------------------------------------------
\maintitle{5. Если дан график производной \(f'(x)\)}

Здесь основное внимание направлено на положение графика относительно
оси \(Ox\).

Если график производной расположен \textbf{выше} оси \(Ox\), то

\[
f'(x)>0,
\]

следовательно, исходная функция возрастает.

Если график производной расположен \textbf{ниже} оси \(Ox\), то

\[
f'(x)<0,
\]

следовательно, исходная функция убывает.

\[
\boxed{
f'(x)>0
\Longrightarrow
f(x)\uparrow
}
\]

\[
\boxed{
f'(x)<0
\Longrightarrow
f(x)\downarrow
}
\]

\subhead{Промежутки знакопостоянства}

Промежутком знакопостоянства называют промежуток, на котором выражение
сохраняет один знак.

Для графика производной это означает:

\begin{itemize}
    \item график целиком выше оси \(Ox\):
    \[
    f'(x)>0;
    \]
    \item график целиком ниже оси \(Ox\):
    \[
    f'(x)<0.
    \]
\end{itemize}

Точки пересечения графика \(f'(x)\) с осью \(Ox\) разбивают ось
на промежутки знакопостоянства.

% -------------------------------------------------
\maintitle{6. Как по \(f'(x)\) находить экстремумы \(f(x)\)}

Сам факт

\[
f'(x_0)=0
\]

ещё не гарантирует экстремум.

Нужно посмотреть, меняется ли знак производной.

\begin{center}
\begin{tabularx}{0.94\linewidth}{@{}c c X@{}}
\toprule
Слева от \(x_0\) & Справа от \(x_0\) & Вывод \\
\midrule
\(f'(x)<0\) & \(f'(x)>0\) &
функция сначала убывает, затем возрастает:
\(x_0\) --- точка локального минимума \\

\(f'(x)>0\) & \(f'(x)<0\) &
функция сначала возрастает, затем убывает:
\(x_0\) --- точка локального максимума \\

\(f'(x)>0\) & \(f'(x)>0\) &
экстремума нет \\

\(f'(x)<0\) & \(f'(x)<0\) &
экстремума нет \\
\bottomrule
\end{tabularx}
\end{center}

Короткая памятка:

\[
-\ \longrightarrow\ +
\qquad
\text{минимум},
\]

\[
+\ \longrightarrow\ -
\qquad
\text{максимум}.
\]

% =================================================
% БЛОК-СХЕМА
% =================================================
\clearpage
\thispagestyle{empty}

\begin{center}
{\Large\bfseries\color{darkaccent}
Алгоритм: как читать задачи по графикам \(f(x)\) и \(f'(x)\)
}
\end{center}

\vspace{7mm}

\begin{center}
\begin{tikzpicture}[
    node distance=9mm and 16mm
]

\node[flowstart] (start)
{Начало};

\node[flowaction, below=of start] (read)
{Прочитайте условие и определите,
что требуется найти};

\node[flowdecision, below=11mm of read] (whose)
{Чей график дан?};

\node[
    flowaction,
    below left=16mm and 19mm of whose
] (function)
{Дан график \(f(x)\).\\
Разбейте его на промежутки
возрастания и убывания};

\node[
    flowaction,
    below right=16mm and 19mm of whose
] (derivative)
{Дан график \(f'(x)\).\\
Разбейте его на промежутки,
где график выше и ниже оси \(Ox\)};

\node[
    flowaction,
    below=10mm of function
] (funconvert)
{Переведите монотонность в знак:
\[
f\uparrow\Rightarrow f'>0,
\qquad
f\downarrow\Rightarrow f'<0
\]};

\node[
    flowaction,
    below=10mm of derivative
] (derconvert)
{Переведите знак производной
в монотонность:
\[
f'>0\Rightarrow f\uparrow,
\qquad
f'<0\Rightarrow f\downarrow
\]};

\coordinate[
    below=16mm of $(funconvert)!0.5!(derconvert)$
] (merge);

\node[
    flowaction,
    below=5mm of merge,
    text width=74mm
] (criterion)
{Сопоставьте полученные интервалы
с критерием из условия задачи};

\node[
    flowaction,
    below=9mm of criterion,
    text width=74mm
] (points)
{Проверьте границы интервалов,
нули производной, выколотые точки
и область определения};

\node[
    flowfinal,
    below=9mm of points
] (answer)
{Запишите именно то,
что требуется в условии:
точки, их количество,
интервал или значение};

\node[
    flowstart,
    below=8mm of answer
] (finish)
{Готово};

\draw[flowarrow] (start) -- (read);
\draw[flowarrow] (read) -- (whose);

\draw[flowarrow]
(whose.west)
-| node[pos=0.27,above,font=\small] {\(f(x)\)}
(function.north);

\draw[flowarrow]
(whose.east)
-| node[pos=0.27,above,font=\small] {\(f'(x)\)}
(derivative.north);

\draw[flowarrow] (function) -- (funconvert);
\draw[flowarrow] (derivative) -- (derconvert);

\draw[flowarrow]
(funconvert.south)
|- (merge);

\draw[flowarrow]
(derconvert.south)
|- (merge);

\draw[flowarrow]
(merge)
-- (criterion);

\draw[flowarrow] (criterion) -- (points);
\draw[flowarrow] (points) -- (answer);
\draw[flowarrow] (answer) -- (finish);

\end{tikzpicture}
\end{center}

\clearpage

% =================================================
% ЕГЭ-АЛГОРИТМ
% =================================================
\maintitle{7. Четыре вопроса перед решением задачи ЕГЭ}

Практически каждую задачу этого типа удобно прогонять через четыре
вопроса.

\subhead{Вопрос 1. Чей график дан?}

\[
f(x)
\qquad\text{или}\qquad
f'(x)?
\]

Это определяет способ чтения рисунка.

\subhead{Вопрос 2. Каков критерий отбора?}

Например:

\[
f'(x)<0,
\]

\[
f'(x)>0,
\]

\[
f'(x)=0,
\]

либо требуется найти точки максимума или минимума функции.

\subhead{Вопрос 3. Что должно происходить на данном графике?}

Если дан \(f(x)\) и требуется

\[
f'(x)<0,
\]

то ищем участки \textbf{убывания функции}.

Если дан \(f'(x)\) и требуется

\[
f'(x)<0,
\]

то ищем участки, где график производной находится
\textbf{ниже оси \(Ox\)}.

\subhead{Вопрос 4. Что именно записывать в ответ?}

Проверяйте формулировку:

\begin{itemize}
    \item число точек;
    \item сами значения \(x\);
    \item количество целых значений \(x\);
    \item точку максимума;
    \item точку минимума;
    \item промежуток возрастания или убывания.
\end{itemize}

\keypoint{Очень важный принцип.}
Даже после правильного анализа графика задача ещё не закончена.
Нужно записать ответ именно в той форме, которая требуется условием.

% -------------------------------------------------
\maintitle{8. Два внешне похожих условия}

Рассмотрим критерий

\[
f'(x)<0.
\]

\subhead{Случай A: дан график функции \(f(x)\)}

Требуется установить, где

\[
f'(x)<0.
\]

Перевод:

\[
f'(x)<0
\quad\Longleftrightarrow\quad
f(x)\text{ убывает}.
\]

Значит, на рисунке ищем участки убывания исходной функции.

\subhead{Случай B: дан график производной \(f'(x)\)}

Снова требуется установить, где

\[
f'(x)<0.
\]

Теперь никакого дополнительного перевода через монотонность для отбора
точек делать не требуется.

Ищем участки, где сам график

\[
y=f'(x)
\]

находится ниже оси \(Ox\).

\begin{center}
\begin{tabularx}{0.94\linewidth}{@{}X X@{}}
\toprule
Дан \(f(x)\) & Дан \(f'(x)\) \\
\midrule
Смотрим, где функция убывает. &
Смотрим, где график лежит ниже \(Ox\). \\[1mm]

Убывание означает \(f'(x)<0\). &
Положение ниже \(Ox\) непосредственно означает \(f'(x)<0\). \\
\bottomrule
\end{tabularx}
\end{center}

% -------------------------------------------------
\maintitle{9. Границы интервалов и особые точки}

Здесь часто теряются баллы.

\subhead{Если требуется строгое неравенство}

При условии

\[
f'(x)<0
\]

точки, где

\[
f'(x)=0,
\]

в ответ не входят.

Аналогично при условии

\[
f'(x)>0.
\]

\subhead{Выколотая точка}

Если на графике производной точка выколота, то соответствующее значение
\(x\) не принадлежит области определения изображённой функции.

Такое значение нельзя учитывать только потому, что оно находится внутри
подходящего по знаку интервала.

\subhead{Конец изображённого промежутка}

Если точка принадлежит графику, проверяйте её по условию отдельно.

Если она выколота или значение функции в ней не определено,
использовать её нельзя.

\subhead{Целые значения}

Если требуется количество целых \(x\), сначала получите допустимые
интервалы, затем перечислите только целые точки внутри них.

Например,

\[
-3<x<2
\]

даёт

\[
x=-2,-1,0,1,
\]

то есть

\[
4
\]

целых значения.

Если условие имеет вид

\[
-3\le x<2,
\]

добавляется точка

\[
x=-3,
\]

и получается уже \(5\) значений.

% -------------------------------------------------
\maintitle{10. Типичные ошибки}

\begin{enumerate}
    \item \textbf{Сразу искать знак, не выяснив, чей график дан.}

    Это приводит к смешению двух разных алгоритмов.

    \item \textbf{На графике \(f(x)\) принимать положение относительно
    оси \(Ox\) за знак производной.}

    Значение самой функции

    \[
    f(x)
    \]

    и знак её производной

    \[
    f'(x)
    \]

    описывают разные свойства.

    \item \textbf{На графике \(f'(x)\) смотреть, возрастает сама
    производная или убывает.}

    Для определения монотонности \(f(x)\) нужен знак \(f'(x)\),
    то есть положение графика производной относительно оси \(Ox\).

    \item \textbf{Считать каждую точку \(f'(x)=0\) экстремумом.}

    Нужна смена знака производной.

    \item \textbf{Включать ноль производной в условие \(f'(x)<0\).}

    Ноль не является отрицательным числом.

    \item \textbf{Не замечать выколотые точки.}

    Они могут изменить количество допустимых целых значений.

    \item \textbf{Записывать сами точки, когда требуется их количество.}

    Всегда перечитывайте последнюю фразу условия.
\end{enumerate}

% -------------------------------------------------
\maintitle{11. Быстрый чек-лист перед ответом}

Проверьте себя:

\begin{itemize}
    \item Я точно определила, чей график дан?
    \item Если дан \(f(x)\), я смотрю на монотонность?
    \item Если дан \(f'(x)\), я смотрю на положение относительно \(Ox\)?
    \item Я правильно перевела условие
    \[
    f'(x)>0
    \quad\text{или}\quad
    f'(x)<0?
    \]
    \item Я проверила точки, где
    \[
    f'(x)=0?
    \]
    \item Я проверила смену знака, если ищу экстремум?
    \item Я исключила выколотые точки?
    \item Я правильно обработала границы интервалов?
    \item Если нужны целые числа, я перечислила только целые допустимые \(x\)?
    \item Я записываю именно тот формат ответа, который требуется?
\end{itemize}

\vspace{3mm}

\begin{center}
{\large\bfseries\color{darkaccent}
Главный ориентир:
сначала определяем, чей график дан.
}
\end{center}


\end{document}